\section{Les évènements dans les transformations d'objets}
\label{sec:formalisation-evenement-transfo-objets}

Considérons deux G-cartes $G$ et $H$. %
On veut qualifier les modifications d'orbites de type $\orbit{\omega}$ entre $G$ et $H$.
Ces dernières peuvent être créées, supprimées, scindées, fusionnées, inchangées ou modifiées.
On parle dans ce cas d'événements survenus sur les orbites.

\subsection{Création}
\label{sec:formalisation-evenement-transfo-objets-creation}



\begin{definition}[Création d'orbite]
	Soient $G$ et $H$ deux G-cartes de dimension $n \in \mathbb{N}$,
	$\orbit{\omega}$ un type d'orbite de dimension $n$ et $v$ un brin de $H$.

	L'orbite $H\orbit{\omega}(v)$ est \emph{créée} si et seulement si aucun élément de l'orbite
	$\orbit{\omega}$ incidente au brin $v$ dans $H$ n'existait
	dans $G$, \textit{i.e.} $G \cap H\orbit{\omega}(v) = \emptyset$. %
\end{definition}

\begin{figure}[H]
	\centering
	\begin{subfigure}{.45\textwidth}
		\centering
		\includegraphics[width=.6\textwidth]{g-carte-triangulation-creation-G}
		\caption{\label{fig:formalisation-evenements-g-carte-evenement-triangulation-sommet-g}$G$}
	\end{subfigure}
	\begin{subfigure}{.45\textwidth}
		\centering
		\includegraphics[width=.6\textwidth]{g-carte-triangulation-creation-H}
		\caption{\label{fig:formalisation-evenement-g-carte-evenement-triangulation-sommet-h}$H$}
	\end{subfigure}
	\caption{\label{fig:formalisation-evenement-g-carte-evenement-triangulation-sommet}Triangulation dans une G-carte~: création d'un sommet \todo{metttre tous les brins}}
\end{figure}

L'orbite sommet $H\orbit{1,2,3}(a_2)$ est créée à l'issue de la transformation
du graphe $G$ vers $H$
(Figure~\ref{fig:formalisation-evenement-g-carte-evenement-triangulation-sommet}). %
En particulier, aucuns des brins $a_2,b_2\ldots h_2$ de $H$ n'existent dans $G$. %


\subsection{Scission}
\label{sec:formalisation-evenement-transfo-objets-scission}


\begin{definition}[Scission d'orbite]
	Soient $G$ et $H$ deux G-cartes de dimension $n \in \mathbb{N}$, $\orbit{\omega}$ un type d'orbite de dimension $n$ et $v$ un brin de $G$.

	L'orbite $G\orbit{\omega}(v)$ est \emph{scindée} si et seulement si il
	existe deux brins $v_1$ et $v_2$ de $G\orbit{\omega}(v)$ tels que $v_1$ et
	$v_2$ sont aussi des brins de $H$ et qu'ils appartiennent à des orbites
	$\orbit{\omega}$ différentes, \textit{i.e.}
	$H\orbit{\omega}(v_1) \neq H\orbit{\omega}(v_2)$.
\end{definition}

\begin{figure}[H]
	\centering
	\begin{subfigure}{.45\textwidth}
		\centering
		\includegraphics[width=.6\textwidth]{g-carte-triangulation-scission-G}
		\caption{\label{fig:formalisation-evenements-g-carte-evenement-triangulation-face-g}$G$}
	\end{subfigure}
	\begin{subfigure}{.45\textwidth}
		\centering
		\includegraphics[width=.6\textwidth]{g-carte-triangulation-scission-H}
		\caption{\label{fig:formalisation-evenement-g-carte-evenement-triangulation-face-h}$H$}
	\end{subfigure}
	\caption{\label{fig:formalisation-evenement-g-carte-evenement-triangulation-face}Triangulation
		dans une G-carte~: scission d'une face \todo{remettre les brins dans le même ordre que pour l'état de l'art}}
\end{figure}

Reprenons l'exemple figure~\ref{fig:formalisation-evenement-g-carte-evenement-triangulation-sommet} %
La face $G\orbit{0,1,3}(a_0)$ est scindée en quatre faces distinctes~:
$H\orbit{0,1,3}(a_0)$,
$H\orbit{0,1,3}(c_0)$,
$H\orbit{0,1,3}(e_0)$,
$H\orbit{0,1,3}(g_0)$ %


\subsection{Suppression}
\label{sec:formalisation-evenement-transfo-objets-suppression}


La suppression est l'inverse de la création. %

\begin{definition}[Suppression d'orbite]
	Soient $G$ et $H$ deux G-cartes de dimension $n \in \mathbb{N}$,
	$\orbit{\omega}$ un type d'orbite de dimension $n$ et $v$ un brin de $G$.

	L'orbite $G\orbit{\omega}(v)$ est \emph{supprimée} si et seulement si
	aucun élément de l'orbite $\orbit{\omega}$ incidente à $v$ dans $G$  n'existe dans $H$, \textit{i.e.} $H \cap G\orbit{\omega}(v) = \emptyset$. %
\end{definition}

Reprenons l'exemple de la
Figure~\ref{fig:formalisation-evenement-g-carte-evenement-triangulation-sommet},
cette fois de $H$ vers $G$. %
Le sommet $H\orbit{1,2,3}(a_{2})$ est supprimé dans $G$ car aucun des brins
$a_2,b_2\ldots h_2$ dans $H$ n'existe dans $G$. %


\subsection{Fusion}
\label{sec:formalisation-evenement-transfo-objets-fusion}


La fusion est l'inverse de la scission. %

\begin{definition}[Fusion d'orbites]
	Soient $G$ et $H$ deux G-cartes de dimension $n \in \mathbb{N}$,
	$\orbit{\omega}$ un type d'orbite de dimension $n$ et $v$ et $v'$ deux brins de $G$. %

	Les orbites distinctes $G\orbit{\omega}(v)$ et $G\orbit{\omega}(v')$ sont
	\emph{fusionnées} si et seulement si il existe deux brins respectivement
	$v_{1}$ de $G\orbit{\omega}(v)$ et $v_{2}$ de $G\orbit{\omega}(v')$ tel que
	$v_{1}$ et $v_{2}$ sont des brins de $H$ de la même orbite $\orbit{\omega}$, \textit{i.e.}
	$H\orbit{\omega}(v_{1}) = H\orbit{\omega}(v_{2})$. %

\end{definition}

% \begin{figure}[H]
%   \centering
%   \includegraphics[width=.8\linewidth]{{figures/}jerboa-merge-then-now}
%   \caption{Fusion de deux sommets voisins d'un cube}
%   \label{fig:jerboa-merge-then-now}
% \end{figure}

Reprenons l'exemple de la figure~\ref{fig:formalisation-evenement-g-carte-evenement-triangulation-sommet}, toujours de $H$ vers $G$.
Les quatre faces $H\orbit{0,1,3}(a_0)$, $H\orbit{0,1,3}(c_0)$, $H\orbit{0,1,3}(e_0)$ et $H\orbit{0,1,3}(g_0)$ sont fusionnées dans $G$. %
En effet, les brins $a_0$, $c_0, e_0$ et $g_0$ appartiennent à la même face dans $G$. %

% représente la transformation d'un cube par la contraction de l'arête entourée en vert.
% Les deux sommets de $G$ entourés rouge sont fusionnés en un seul sommet dans $H$ entouré en rouge.
% Deux orbites sont considérées comme fusionnées en une seule si leurs nœuds préservés sont, après
% transformation, réunis dans une seule et même orbite.
% Formellement, pour deux nœuds de $G$, $v, v' \in V^2$ d'orbites différentes (c'est-à-dire telles que $G\orbit{\omega}(v) \neq G\orbit{\omega}(v')$), la \emph{fusion} de deux orbites est définie de la façon suivante~:

% Ces deux orbites $G\orbit{\omega}(v)$ et $G\orbit{\omega}(v')$
% sont \emph{fusionnées} s'il existe deux nœuds $v_1$ et $v_2$ appartenant respectivement à $G\orbit{\omega}(v)$ et $G\orbit{\omega}(v')$ et tous les deux nœuds de $H$, tels que $H\orbit{\omega}(v_1) = H\orbit{\omega}(v_2)$.


\subsection{Non-changement}
\label{sec:formalisation-evenement-transfo-objets-non-changement}


\begin{definition}[Non-changement d'orbite]
	Soient $G$ et $H$ deux G-cartes de dimension $n \in \mathbb{N}$,
	$\orbit{\omega}$ un type d'orbite de dimension $n$ et $v$ un brin de $G$.

	L'orbite $G\orbit{\omega}(v)$ est \emph{inchangée} si et seulement si elle identique à $H\orbit{\omega}(v)$, \textit{i.e.} $G\orbit{\omega}(v) = H\orbit{\omega}(v)$.
\end{definition}

Dans la figure~\ref{fig:formalisation-evenement-g-carte-evenement-triangulation-sommet}, l'arête incidente au brin $a_0$ est inchangée, puisqu'elle est identique dans $G$ et $H$. %

% \begin{figure}[H]
%   \centering
%   \includegraphics[width=.8\linewidth]{{figures/}jerboa-unchange-example}
%   \caption{Sommets inchangés par l'insertion d'un sommet sur une arête}
%   \label{fig:jerboa-unchange-example}
% \end{figure}

% Dans l'exemple de la figure \textbf{fig:jerboa-unchange-example}, l'insertion d'un
% sommet sur une arête \emph{ne change pas} les sommets aux extrémités de cette
% arête (encerclés en rouge). Une orbite est considérée comme \emph{inchangée} si
% elle est préservée, c'est-à-dire qu'elle est identique dans $G$ et $H$.
% Formellement, pour nœud de $G$, $v \in V$, le \emph{non-changement} d'une orbite
% est défini de la façon suivante~:

% Une orbite $G\orbit{\omega}(v)$ est \emph{inchangée} si l'orbite $G\orbit{\omega}(v)$ est aussi une orbite de $H$, c'est-à-dire si $G\orbit{\omega}(v) \subset H$.


\subsection{Modification}
\label{sec:formalisation-evenement-transfo-objets-modification}


\begin{definition}[Modification d'orbite]
	Soient $G$ et $H$ deux G-cartes de dimension $n \in \mathbb{N}$,
	$\orbit{\omega}$ un type d'orbite de dimension $n$ et $v$ un brin de $G$. %

	L'orbite $G\orbit{\omega}(v)$ est \emph{modifiée} si et seulement si il existe un brin $v'$ de
	$G\orbit{\omega}(v)$ et $H$ tel que tout brin préservé (de $G \cap H$) $v''$ appartient à $G\orbit{\omega}(v')$ si et seulement si $v''$ appartient à $H\orbit{\omega}(v')$%
\end{definition}

Dans la
figure~\ref{fig:formalisation-evenement-g-carte-evenement-triangulation-face},
le sommet $G\orbit{1,2,3}(a_{0})$ est modifié dans $H$, sans fusion ni scission, puisque
$G\orbit{1,2,3}(a_{0}) \neq H\orbit{1,2,3}(a_{0})$. %

% \maxime{Ajouter l'exemple d'insertion de sommet en entourant la face modifiée}
% La figure ? reprend l'exemple de l'insertion d'un sommet sur une arête du cube. La face entourée en rouge dans $G$ est modifiée dans $H$, car tous ses brins préservés appartiennent à la même face entourée en rouge. Cette face contient des brins supplémentaires du sommet ajouté.\\

% \maxime{Ajouter l'exemple de contraction d'arête en entourant la face modifiée}
% La figure ? reprend l'exemple de la contraction d'arête d'un cube. La face entourée en rouge dans $G$ est modifiée dans $H$, car tous ses brins préservés appartiennent à la même face entourée en rouge.
% Cette face ne contient plus les brins de l'arête contractée.\\

% \maxime{Ajouter l'exemple d'un edgeflip sur une face du cube en entourant tout le volume}
% La figure ? représente le \textit{edge flipping} de l'arête entourée en vert. Le volume entouré en rouge est modifié, car si il contient les mêmes brins dans $G$ et dans $H$, ses liaisons sont modifiées.
% Formellement, pour un nœud de $G$, $v \in V$. La \emph{modification} d'une orbite est définie de la façon suivante~:

% $G\orbit{\omega}(v)$ est modifiée dans $H$ si il existe un nœud $v'$ de $G\orbit{\omega}(v)$ et de $H$ tel que tout nœud préservé $v''$ de $G\orbit{\omega}(v')$ appartient $H\orbit{\omega}(v')$ et réciproquement tout nœud préservé $v''$ de $H\orbit{\omega}(v')$ appartient à $G\orbit{\omega}(v')$.

%En particulier si l'orbite $G\orbit{\omega}(v)$ possède un nœud conservé $v$, alors $G\orbit{\omega}(v)$ ne peut être une orbite supprimée, et $H\orbit{\omega}(v)$ ne peut être une orbite créée.

%%% Local Variables:
%%% mode: latex
%%% TeX-master: "../../../main"
%%% End:
